\documentclass[11pt,aspectratio=169]{ctexbeamer}

\usetheme{default}
\setbeamertemplate{navigation symbols}{}

\usepackage{amsmath}
\usepackage{booktabs}

\definecolor{DeepBlue}{RGB}{25,74,122}
\definecolor{SoftInk}{RGB}{53,59,66}
\definecolor{PaperBg}{RGB}{249,248,245}
\definecolor{PanelBlue}{RGB}{236,243,249}

\setbeamercolor{normal text}{fg=SoftInk,bg=PaperBg}
\setbeamercolor{structure}{fg=DeepBlue}
\setbeamercolor{frametitle}{fg=DeepBlue,bg=PaperBg}
\setbeamercolor{title}{fg=DeepBlue}
\setbeamercolor{block title}{fg=DeepBlue,bg=PanelBlue}
\setbeamercolor{block body}{fg=SoftInk,bg=white}

\newcommand{\Core}[1]{
  \begin{block}{核心结论}
    #1
  \end{block}
}

\newcommand{\LectureMeta}[2]{
  \title{大模型推理基础设施}
  \subtitle{#1}
  \author{张书豪}
  \institute{华中科技大学计算机学院}
  \date{#2}
}


\LectureMeta{第 5 讲 KV Cache 管理}{研究生课程课件}

\begin{document}

\begin{frame}[plain]
  \titlepage
\end{frame}

\begin{frame}{教学目标}
  \Core{本讲将 KV Cache 作为影响并发、调度与服务稳定性的核心系统对象进行讨论。}
  \begin{itemize}
    \item 理解 KV 生命周期与容量压力。
    \item 理解分页、块管理、碎片与回收。
    \item 说明 KV 组织方式如何反向影响调度器能力。
  \end{itemize}
\end{frame}

\begin{frame}{为什么 KV 是系统核心对象}
  \Core{KV 既是性能关键状态，又是容量主约束，因此它会塑造系统设计空间。}
  \begin{itemize}
    \item 随上下文长度和并发增长快速膨胀。
    \item 需要跨 decode 多轮驻留。
    \item 既影响时延，也限制 batch 扩展空间。
  \end{itemize}
\end{frame}

\begin{frame}{为什么 KV 问题会在大模型时代被放大}
  \Core{随着上下文更长、并发更高、交互更持续，KV 已经从实现细节变成系统一等约束。}
  \begin{itemize}
    \item 长上下文让单请求状态膨胀。
    \item 在线服务让状态必须跨轮驻留。
    \item 多租户与动态到达让分配与回收更频繁。
  \end{itemize}
\end{frame}

\begin{frame}{KV 生命周期}
  \Core{理解 KV，首先要理解它如何生成、追加、复用、迁移和回收。}
  \begin{enumerate}
    \item prefill 生成首轮 KV
    \item decode 持续追加
    \item 可能发生前缀复用或跨请求共享
    \item 完成后释放或转入更低层存储
  \end{enumerate}
\end{frame}

\begin{frame}{连续地址为什么会成为问题}
  \Core{如果逻辑连续序列必须映射到物理连续大块显存，动态服务很快会受到碎片和预留浪费限制。}
  \begin{itemize}
    \item 并发请求越多，连续空间越难保证。
    \item 预留过大则浪费，预留过小则频繁失败。
    \item 这正是分页式 KV 管理出现的直接背景。
  \end{itemize}
\end{frame}

\begin{frame}{PagedAttention 的设计直觉}
  \Core{PagedAttention 的关键在于将逻辑连续状态映射为物理离散块。}
  \begin{itemize}
    \item 避免显存必须连续分配大块空间。
    \item 降低内部碎片与预留浪费。
    \item 为更动态的请求组织提供底层支撑。
  \end{itemize}
\end{frame}

\begin{frame}{PagedAttention 的系统意义}
  \Core{PagedAttention 的价值，不只在于节省显存，更在于把状态组织方式改造成了支持动态调度的形式。}
  \begin{itemize}
    \item 更灵活的状态分配与回收。
    \item 更小的显存预留压力。
    \item 更强的 continuous batching 支撑能力。
  \end{itemize}
\end{frame}

\begin{frame}{为什么 KV 会反向约束调度}
  \Core{调度器看似管理请求，但真正决定并发上限的常常是状态驻留能力。}
  \begin{itemize}
    \item 如果 KV 不能灵活管理，continuous batching 只会更快碰到容量墙。
    \item 如果回收和复用代价过高，调度收益会被状态管理开销抵消。
    \item 这就是为什么 vLLM 的调度贡献必须和 PagedAttention 一起理解。
  \end{itemize}
\end{frame}

\begin{frame}{前缀复用为何进一步提升复杂度}
  \Core{一旦系统开始尝试前缀复用，KV 就不再只是“每个请求私有状态”，而开始成为跨请求共享对象。}
  \begin{itemize}
    \item 共享带来更高潜在收益。
    \item 也带来正确性、元数据组织和回收边界问题。
    \item 这使 KV 进一步接近广义状态管理问题。
  \end{itemize}
\end{frame}

\begin{frame}{源码穿插：nano-vllm-hust 中的 KV 管理结构}
  \Core{\texttt{\detokenize{nanovllm/engine/block_manager.py}} 将课堂中的块管理、引用计数与前缀哈希机制压缩到了一个较短文件中。}
  \begin{itemize}
    \item \texttt{\detokenize{Block}} 维护块编号、引用计数、哈希值与块内 token。
    \item \texttt{\detokenize{free_block_ids}} 与 \texttt{\detokenize{used_block_ids}} 体现了显式的块生命周期管理。
    \item \texttt{\detokenize{hash_to_block_id}} 为前缀块复用提供查找入口。
  \end{itemize}
\end{frame}

\begin{frame}{源码穿插：课堂阅读建议}
  \Core{阅读 \texttt{\detokenize{BlockManager}} 时，建议按“能否分配、如何分配、何时回收、何时哈希”四个问题组织。}
  \begin{itemize}
    \item \texttt{\detokenize{can_allocate()}} 体现了前缀缓存命中与新块需求估算。
    \item \texttt{\detokenize{allocate()}} 和 \texttt{\detokenize{deallocate()}} 对应块表构造与引用计数回收。
    \item \texttt{\detokenize{hash_blocks()}} 则把完整块写回全局哈希表，为后续复用建立条件。
  \end{itemize}
\end{frame}

\begin{frame}{例子：算一遍 KV 容量为什么能立刻改变系统判断}
  \Core{KV 课程里最应该让学生亲手做的，是容量算术，因为它直接决定你如何理解并发上限。}
  \begin{block}{输入更长}
    单请求 KV 规模线性增长，prefill 结束后就把一大块状态留在系统里。
  \end{block}
  \pause
  \begin{block}{并发更高}
    多请求同时驻留时，容量压力不是加法直觉那么简单，还会叠加块管理与碎片问题。
  \end{block}
  \pause
  \begin{block}{课堂结论}
    只要状态规模算错，后续对调度、throughput 和 tail latency 的解释就会整体偏掉。
  \end{block}
\end{frame}

\begin{frame}{例子：为什么连续大块分配很快会失败}
  \Core{传统“给每个请求预留一整段连续显存”的想法，在动态服务下几乎注定碰壁。}
  \begin{block}{预留过大}
    会浪费大量空闲容量，使理论并发远高于实际并发。
  \end{block}
  \pause
  \begin{block}{预留过小}
    一旦上下文稍长，请求就会频繁扩容甚至失败。
  \end{block}
  \pause
  \begin{block}{PagedAttention 的意义}
    把逻辑连续变成物理分块后，系统终于可以更诚实地管理状态生命周期。
  \end{block}
\end{frame}

\begin{frame}{代码例子：BlockManager 里最值得学生先看的对象}
  \Core{KV 讲义里的代码例子，应该从对象关系入手，而不是从论文术语入手。}
  \begin{block}{核心对象}
    {\ttfamily\small
    \detokenize{Block(block_id, ref_count, hash_value, token_ids)}\par
    \detokenize{free_block_ids}\par
    \detokenize{used_block_ids}\par
    \detokenize{hash_to_block_id}\par
    }
  \end{block}
  \pause
  \begin{block}{课堂应追问}
    这些对象分别体现了分配、占用、共享和回收的哪一面？
  \end{block}
  \pause
  \begin{block}{系统含义}
    读懂这些对象，才会知道 KV 管理讲的并不是“抽象缓存”，而是明确的数据结构与状态机。
  \end{block}
\end{frame}

\begin{frame}{代码例子：从 can\_allocate 到 allocate 的判断链}
  \Core{课堂最适合拿来解释的，不是整份代码，而是“能否分配 -> 如何分配”的判断链。}
  \begin{block}{代码片段}
    {\ttfamily\small
    \detokenize{num_required_blocks = seq.num_blocks - num_computed_blocks}\par
    \detokenize{return num_required_blocks <= len(self.free_block_ids)}\par
    \detokenize{seq.block_table.append(block_id)}\par
    }
  \end{block}
  \pause
  \begin{block}{课堂应追问}
    为什么“还剩多少 free blocks”本身就已经在决定系统能否扩大并发？
  \end{block}
  \pause
  \begin{block}{结论}
    KV 不是调度的附属品，容量判断本身就在塑造调度器的可选动作空间。
  \end{block}
\end{frame}

\begin{frame}{例子：前缀复用为什么比普通缓存更复杂}
  \Core{一旦系统尝试前缀共享，KV 就从“每请求私有状态”变成“可能跨请求共享状态”。}
  \begin{block}{收益直觉}
    相同前缀意味着重复 prefill 可被避免，潜在节省显著。
  \end{block}
  \pause
  \begin{block}{复杂性来源}
    共享需要元数据索引、引用计数和正确回收边界，否则收益可能变成一致性问题。
  \end{block}
  \pause
  \begin{block}{课程边界}
    课堂里可以先讲 observe-only 的潜在命中，再逐步进入真实复用机制。
  \end{block}
\end{frame}

\begin{frame}{课堂练习：把 KV 问题翻译成系统因果链}
  \Core{KV 这一讲不应该只留下术语，而应该留下至少一条学生自己能复述的因果链。}
  \begin{block}{练习一}
    用一句因果链解释：为什么长上下文会通过状态驻留拖慢短请求体验？
  \end{block}
  \pause
  \begin{block}{练习二}
    说明为什么分页/块管理的系统意义不只是“节省显存”。
  \end{block}
  \pause
  \begin{block}{练习三}
    如果你要为 BlockManager 加第一条观测日志，你会放在 can\_allocate、allocate 还是 deallocate？为什么？
  \end{block}
\end{frame}


\begin{frame}{课堂讨论}
  \Core{KV 管理不仅涉及显存节约，更直接关系到系统可扩展性。}
  \begin{itemize}
    \item 为什么 KV 不应被视为 attention 的临时中间量？
    \item 为什么分块管理可以扩大系统控制空间？
    \item 当前缀复用进入系统后，正确性与收益之间有什么张力？
  \end{itemize}
\end{frame}

\begin{frame}{本讲总结}
  \Core{理解 KV 管理机制，是分析 vLLM 一类系统如何将调度收益扩展到更高并发的前提。}
  \begin{itemize}
    \item KV 是容量问题，也是调度问题。
    \item 后续状态统一视角会继续扩展这一点。
  \end{itemize}
\end{frame}

\end{document}